\section{Introduction}

An odd cycle in the permutation induced by two lines of a Latin square
supports a simple trade: swap the two lines on that cycle.  The move preserves
the Latin property.  In the row and column views it reverses the conventional
row--column Latin sign; in the symbol view it preserves that sign and reverses
the two parastrophically transported signs. Wanless already gives the
underlying three line-parity changes for cycle
switches~\cite[Proposition~1]{Wanless2004CycleSwitches}; the pairwise sign
products here are their explicit corollary. This local operation is therefore
a natural primitive in parity-based approaches related to the Alon--Tarsi
setting, but its three sign conventions must be kept distinct.  B\'erczi's
experimental program studies such local corrections as ingredients for
global constructions~\cite{Berczi2026}; it does not supply the theorems proved
here.

We isolate the underlying existence question.  For each of the row, column,
and symbol views, ask whether some pair of lines induces a permutation with a
nontrivial odd cycle.  The three answers form a pattern in $\{F,T\}^3$.  A
square is FFF when all three answers are negative.  Equivalently, every union
of two link factors in the tripartite representation has component lengths
divisible by four. The two-factor/cycle correspondence in one view is
classical~\cite[Section~3]{WanlessIhrig2005Factorizations}; the three-view
formulation follows by conjugating the coordinate roles.

The structural results start with two elementary boundaries.  Every
odd-order Latin square is TTT, whereas the Cayley table of a finite group
is FFF exactly when the group is a
$2$-group.  A basis-family closure criterion detects group isotopy.  We then
show that failure patterns are monotone from Latin subsquares and analyze
quasigroup congruences: odd block size forces TTT, while two-element blocks
preserve the quotient pattern exactly.  Thus a hypothetical FFF quasigroup of
order $2p$, for odd prime $p$, must be congruence-simple in every isotope.
This excludes all congruence-imprimitive constructions at order $10$ but does
not by itself decide the simple case.

The remaining binary line-incidence defect also has an exact subsquare
meaning.  Nonuniversal vectors in the left binary line kernel are in
bijection with half-order subsquares, and the quotient classes group them
into complementary four-packs.  This yields a closed counting formula and
forces maximal binary line rank for every FFF order congruent to $2$ modulo
$4$ and greater than $2$.  At order $10$ this is necessary, not sufficient.

More generally, Moorhouse's $3$-net character formula specializes as follows:
a nonuniversal line dependency over a prime field is a
balanced cyclic quotient homotopy, and after isotopy its fibres form a
quasigroup congruence.  Consequently every Smith torsion prime divides the
order.  For order $2p$, with $p$ odd prime, the congruence obstruction shows
that even one F view forces a saturated integer line-incidence lattice.

The group construction is not the whole phenomenon.  A twisted binary
extension of $C_4$ gives an explicit non-group-isotopic FFF loop of order $8$.
Moreover, direct products satisfy the exact componentwise formula
\[
  \pat(L\times K)=\pat(L)\mathbin{\mathrm{OR}}\pat(K),
\]
and a product is group-isotopic exactly when both factors are.  These theorems
turn one order-$8$ example into infinite power-of-two families.

The order-$8$ computation scans the cited main-class census.  McKay,
Meynert, and Myrvold enumerated the main classes of small Latin
squares~\cite{McKayMeynertMyrvold2007};
the associated ANU file contains one representative of each of the $283657$
order-$8$ main classes~\cite{McKayLatinData}.  The scan finds exactly $230$
FFF main classes.  Only $5$ are group-isotopic and $225$ are not, and all
eight ordered failure patterns occur.  Independently, fresh exact-cover
generation of all $9408$ reduced order-$6$ squares finds no FFF square and
matches the official ANU set exactly.

Combining theory and computation proves that every order $2^m$, $m\geq3$,
realizes all eight patterns and contains a non-group-isotopic FFF square.  We
do not claim unrestricted direct-product cancellation on isotopy or main
classes. Section~\ref{sec:fibred-growth} establishes more specific
distinction and classification results under explicit hypotheses.
The accompanying
complete research report proves by exhaustive master searches that order $10$
has no FFF square; its proof is not repeated in this compact companion,
but its result is explicitly used in the order-$20$ quotient-freeness
argument and in $h(5)=2$. A reduced order-$12$ FFF witness, independently checked in all three
views, refutes the former power-of-two conjecture; orders $14$ and $36$ have
further positive witnesses. The classical prime round-robin table is FFF
exactly for $p\equiv3\pmod8$, supplying infinitely many non-power-of-two
orders; with the separately audited order-$10$ exclusion, its order-$20$
instance gives $h(5)=2$. A written three-point lift proof gives an FFF
square of order $m2^{\operatorname{ord}_m(2)}$ for every odd $m>1$, without
determining the least dyadic exponent in general.  Order $18$ remains the
next undecided small even non-power-of-two order.

The row-fibred extension in Section~\ref{sec:fibred-growth} generalizes
the product pattern formula. Quotient-free fibres make the quotient layers
intrinsic; over an elementary abelian outer group this gives an exact
affine-orbit classification within a fixed family. Exact finite controls
supply the $514$ order-$20$ input classes. The resulting large counts
come from a construction theorem and Burnside's lemma, not direct scans
of every output square.

The longer companion develops additional cycle-parity, subset-action,
primitive-group, and flag-surface machinery, gives the order-$10$
completeness argument, and records the restricted order-$18$ search
boundary. The general product, row-fibred, and quotient-descent
proofs do not depend on that machinery; their stated order-$20$
specialization does depend on the order-$10$ computation.
